% ============================================================================
% CAPÍTULO 2 — REFERENCIAL TEÓRICO
% ============================================================================
\chapter{Referencial Teórico}
\label{cap:referencial}

\section{Personalidade e o modelo Big Five}

\subsection{O que é personalidade}

A personalidade tem sido conceituada a partir de diversas perspectivas
teóricas e em diferentes níveis de abstração, e um dos níveis mais
estudados é o dos traços de personalidade \cite{johnsrivastava1999}. A
abordagem de traços, adotada neste trabalho, parte do princípio de que as
diferenças individuais de personalidade podem ser descritas por um número
pequeno de dimensões amplas e relativamente estáveis ao longo do tempo e
das situações, chamadas \textit{traços de personalidade}. Um traço não é
uma categoria discreta (a pessoa ``é'' ou ``não é'' extrovertida), mas um
\textit{continuum}: cada indivíduo pode ser posicionado em algum ponto
entre dois polos de cada dimensão.

\subsection{O modelo Big Five}

O modelo Big Five (também chamado de modelo dos cinco grandes fatores, ou
\textit{Five-Factor Model}) é a taxonomia de traços de personalidade em
torno da qual a área tem se aproximado de um consenso
\cite{johnsrivastava1999}. Ele emergiu da convergência de duas tradições
de pesquisa independentes: a abordagem lexical, que parte da análise
fatorial de adjetivos de personalidade usados na linguagem natural
\cite{goldberg1990}, e a abordagem baseada em questionários
psicométricos, como o \textit{Revised NEO Personality Inventory}
(NEO-PI-R) \cite{costamccrae1992}. Ambas as tradições chegaram a uma
estrutura de cinco dimensões amplas e relativamente independentes entre
si \cite{johnsrivastava1999}:

\begin{itemize}
    \item \textbf{Extroversão}: tendência a ser sociável, assertivo,
    ativo e a buscar estimulação e a companhia de outras pessoas, em
    oposição a um estilo mais reservado e independente;
    \item \textbf{Amabilidade} (também chamada de Afabilidade): tendência
    a ser cooperativo, confiante e altruísta em relação aos outros, em
    oposição a um estilo mais cético, competitivo ou antagônico;
    \item \textbf{Consciência}: tendência a ser organizado,
    disciplinado, orientado a objetivos e cumpridor de obrigações, em
    oposição a um estilo mais despreocupado ou impulsivo;
    \item \textbf{Neuroticismo}: tendência a experimentar emoções
    negativas (ansiedade, insegurança, instabilidade emocional) com maior
    frequência e intensidade, em oposição à estabilidade emocional;
    \item \textbf{Abertura à Experiência}: tendência à curiosidade
    intelectual, à imaginação e à apreciação por ideias, experiências
    estéticas e situações não convencionais, em oposição a um estilo mais
    convencional e pragmático.
\end{itemize}

Diversos instrumentos psicométricos foram desenvolvidos para medir esses
cinco traços, com diferentes números de itens e propósitos. O estudo
piloto deste trabalho utilizou o \textbf{BFI-10} \cite{rammstedtjohn2007},
uma versão reduzida de 10 itens (2 por traço) do \textit{Big Five
Inventory}, na tradução para o português utilizada por
\citeonline{debortoli2019} (detalhes no Capítulo \ref{cap:metodologia}). A
análise principal deste trabalho utiliza o \textbf{BFI-44}
\cite{johndonahuekentle1991}, versão completa de 44 itens, em uma versão em
português que segue a ordem e o conteúdo dos itens do instrumento
original; a validação psicométrica do BFI-44 para o Brasil foi feita por
\citeonline{andrade2008} (detalhes na seção \ref{sec:bfi44}). Na
literatura revisada neste capítulo aparece ainda o SAPA-20, a forma de 20
itens do \textit{SAPA Personality Inventory} \cite{condonrevelle2015},
utilizada por \citeonline{jiangpengyan2024}.

\section{Finanças comportamentais e perfil de investidor}

O \textbf{perfil de investidor} é o resultado de um processo de avaliação,
em geral um questionário estruturado, que busca estimar o quanto um
indivíduo está disposto a assumir risco em suas decisões de investimento
(tolerância) e é capaz de fazê-lo (capacidade financeira), classificando-o
tipicamente em categorias como conservador, moderado e arrojado. No
mercado financeiro brasileiro, essa avaliação, conhecida como
\textit{suitability}, é uma exigência regulatória: instituições que
distribuem produtos de investimento devem verificar a adequação desses
produtos ao perfil do cliente antes de recomendá-los
\cite{cvm30_2021}. O instrumento utilizado neste trabalho, o IPA (Índice
de Perfil do Investidor), segue esse modelo: é uma adaptação do
questionário de análise de perfil do investidor do Banco do Brasil,
alinhado às diretrizes da Associação Brasileira das Entidades dos
Mercados Financeiro e de Capitais (ANBIMA) e da \textit{International
Organization of Securities Commissions} (IOSCO), com oito perguntas
objetivas sobre horizonte de investimento, patrimônio, experiência prévia
e reação a cenários de perda (detalhes no Capítulo \ref{cap:metodologia}).

Uma linha de crítica relevante a esses instrumentos, levantada por
\citeonline{debortoli2019}, é que questionários de perfil de investidor
como o IPA, baseados em autorrelato sobre preferências e comportamentos
financeiros declarados, podem não capturar integralmente o comportamento
de risco \textit{real} de um indivíduo, medido, por exemplo, por meio de
simulações de investimento com dados reais de rentabilidade. No estudo
dos autores, o IPA teve relação coerente com o comportamento observado em
uma simulação de investimento, mas não foi suficiente, isoladamente, para
explicar todo esse comportamento, resultado que os autores interpretam
como evidência de que fatores adicionais, entre eles traços de
personalidade, contribuem para explicar o comportamento de risco além do
que o questionário de perfil de investidor capta sozinho. É essa lacuna,
a relação entre personalidade e perfil de investidor, que motiva a linha
de pesquisa revisada na seção seguinte.

\section{Personalidade e decisões de investimento: evidências empíricas}

\subsection{Jiang, Peng e Yan (2024)}
\label{sec:jiang}

\citeonline{jiangpengyan2024} investigaram a relação entre os cinco
traços do Big Five e o comportamento de investimento em uma amostra de
3.325 investidores americanos, associados da \textit{American Association
of Individual Investors} (AAII). Os traços de personalidade foram medidos
pelo SAPA-20 \cite{condonrevelle2015}, um instrumento de 20 itens (4 por
traço, escala de 1 a 6), e relacionados, por meio de regressões OLS com
os cinco traços como regressores simultâneos e controle por variáveis
demográficas, a quatro grupos de variáveis: crenças sobre o mercado
(retorno esperado, probabilidade de \textit{crash}, crescimento do
produto interno bruto, PIB, e inflação esperada), aversão ao risco,
comportamento de manada (\textit{herding}) e alocação de carteira
(percentual do patrimônio aplicado em ações).

Os principais achados dos autores são os seguintes. Primeiro,
\textbf{Neuroticismo} e \textbf{Abertura à Experiência} são os dois traços
que os próprios autores destacam como de maior poder explicativo sobre as
decisões de investimento; isso não significa que os outros três traços
tenham todos papel secundário ou nulo (ver o terceiro achado). Segundo,
os autores associam a cada um desses dois traços um canal distinto sobre
a alocação de carteira. No caso do Neuroticismo, o canal provável é o das
crenças: quem pontua mais alto nesse traço tende a avaliar o mercado de
forma mais pessimista, atribuindo maior probabilidade a um \textit{crash},
esperando menor crescimento do PIB e inflação mais alta, e aloca uma
proporção menor do patrimônio em ações. No caso da Abertura à
Experiência, o canal principal é a preferência por risco: Abertura mais
baixa está associada a maior aversão ao risco e a uma alocação menor em
ações, sem relação com o retorno médio esperado do mercado. Os autores
ressalvam, porém, que os dois traços continuam significativos mesmo
controlando a aversão ao risco e as expectativas de retorno, o que indica
a presença também de outros canais. Terceiro, Consciência tem
um padrão próprio, menos enfatizado pelos autores no resumo do artigo:
associa-se a expectativas \textit{mais otimistas} sobre o retorno
esperado do mercado e, ao mesmo tempo, a um investimento menor
em ações, efeito significativo na amostra americana e na alemã (nesta,
medido pela participação no mercado de ações), mas não na australiana,
em que só atinge significância ao nível de 10\% em uma das duas
especificações. Ou seja, Consciência tem efeito sobre a alocação de
carteira em duas das três amostras, na direção oposta à sugerida pelas
expectativas mais otimistas associadas a esse traço. Já Extroversão
associa-se a expectativas mais otimistas sobre o mercado (retorno
esperado e crescimento do PIB) e a menor aversão ao risco declarada, mas,
ao contrário de Consciência, sem efeito significativo sobre a
alocação de carteira nas amostras americana e australiana; apenas na
amostra alemã, em que a variável é a participação no mercado de ações,
Extroversão aparece associada negativamente a essa participação. Quanto
à Amabilidade, o traço está associado a maior aversão ao risco declarada
nas perguntas hipotéticas da pesquisa, mas sem reflexo na alocação real
de carteira; os autores não encontram evidência de que ele tenha papel
direto nas decisões financeiras. Por último, os autores
relatam que Neuroticismo e Extroversão também se associam a maior
comportamento de manada, e que o poder explicativo conjunto dos cinco
traços sobre as crenças do investidor é comparável ao de todas as
variáveis demográficas combinadas (gênero, idade, renda, riqueza,
educação e localização), resultado que os autores replicaram em uma
amostra de investidores chineses. Já as associações de Neuroticismo e
Abertura à Experiência com o investimento em ações foram replicadas em
amostras representativas da Austrália e da Alemanha.

\subsection{De Bortoli, da Costa Jr., Goulart e Campara (2019)}

\citeonline{debortoli2019} compararam quatro paradigmas teóricos
distintos, a saber: a Teoria da Perspectiva, o questionário de perfil de
investidor IPA (Banco do Brasil/ANBIMA), o BFI-10 (Big Five) e o
\textit{Cognitive Reflection Test} (CRT), um teste de reflexão
cognitiva, quanto à sua capacidade de explicar o perfil de risco
\textit{real} de 124 estudantes de graduação da Universidade Federal de
Santa Catarina (UFSC). Esse perfil de risco real foi medido por meio de
uma simulação computacional de investimento (\textit{ExpEcon}, 18
períodos), com ativos oferecidos pelo Banco do Brasil e dados históricos
reais de rentabilidade, e não apenas por autorrelato. A participação foi
voluntária e sem nenhum tipo de incentivo material aos participantes.

Para o Big Five, os autores utilizaram o BFI-10 e categorizaram o escore
de cada traço em três níveis (Baixo, Moderado e Alto), em vez de
tratá-lo como uma variável contínua. Essa categorização é a base direta
do desenho experimental adotado neste trabalho, que usa três níveis
equivalentes (Baixo, Médio e Alto), com a diferença de que aqui os
níveis são induzidos, e não medidos (Capítulo \ref{cap:metodologia}).
Usando regressão logística ordenada (variável dependente: o perfil de
risco obtido no simulador, em três categorias), os autores encontraram
que, de todos os cinco traços do Big Five, apenas Abertura à Experiência
foi estatisticamente significativa: maior
Abertura esteve associada a maior chance de o participante apresentar um
perfil ``arrojado'' no simulador. Os outros quatro traços não
apresentaram relação significativa com o perfil de risco real.

Os autores relatam ainda dois achados adicionais relevantes para este
trabalho. Primeiro, como já mencionado na seção anterior, o IPA não
explicou isoladamente o comportamento de risco real dos participantes.
Segundo, no questionário de Teoria da Perspectiva aplicado antes da
simulação, 92\% dos participantes responderam predominantemente de forma
incompatível com a Teoria da Utilidade Esperada, um resultado consistente
com a Teoria da Perspectiva e que reforça, de forma mais ampla, a relevância de fatores
psicológicos, e não apenas racionais e econômicos, na explicação do
comportamento de investimento.

\subsection{Síntese e hipóteses de referência}
\label{sec:hipoteses-referencia}

Os dois estudos revisados nesta seção convergem parcialmente, mas não
totalmente, quanto a quais traços do Big Five são relevantes para o
comportamento de investimento, o que já é, em si, um indício de que essa
relação pode depender do contexto, da amostra e de como o comportamento de
risco é medido: por autorrelato de alocação de carteira, em
\citeonline{jiangpengyan2024}, e por simulação de investimento com dados
históricos reais, em \citeonline{debortoli2019}. Ainda assim,
quatro pontos de comparação emergem como \textbf{hipóteses de
referência} para os resultados deste trabalho (Capítulo
\ref{cap:discussao}):

\begin{itemize}
    \item \textbf{Abertura à Experiência} é o traço mais consistentemente
    associado à tolerância ao risco de investimento nos dois estudos:
    significativo tanto em \citeonline{debortoli2019} (o único traço
    significativo do estudo) quanto em \citeonline{jiangpengyan2024} (um
    dos dois traços com maior poder explicativo);
    \item \textbf{Neuroticismo} tem papel relevante em
    \citeonline{jiangpengyan2024} (via crenças pessimistas sobre o
    mercado), mas não foi significativo em \citeonline{debortoli2019}, uma
    divergência que pode refletir diferenças de amostra (investidores
    reais vs. estudantes) ou de variável dependente (crenças e alocação
    declarada vs. comportamento em simulação);
    \item \textbf{Consciência} segue o mesmo padrão de divergência que o
    Neuroticismo: tem efeito negativo sobre a alocação de carteira em
    \citeonline{jiangpengyan2024} (em duas das três amostras do artigo),
    mas não foi significativa em \citeonline{debortoli2019};
    \item \textbf{Amabilidade} não teve relação significativa em
    \citeonline{debortoli2019} e, em \citeonline{jiangpengyan2024},
    associou-se a maior aversão ao risco declarada, sem efeito sobre a
    alocação de carteira.
\end{itemize}

A categorização dos níveis de personalidade em Baixo/Moderado/Alto
adotada por \citeonline{debortoli2019} e a regressão OLS com os cinco
traços como regressores simultâneos empregada por
\citeonline{jiangpengyan2024} formam, em conjunto, a base metodológica do
desenho experimental adotado neste trabalho, detalhado no Capítulo
\ref{cap:metodologia}.

\section{Agentes de IA e modelos de linguagem como sujeitos simulados}

A pesquisa revisada na seção anterior foi conduzida com sujeitos humanos,
e nenhum dos dois artigos-base trata do uso de modelos de linguagem como
respondentes. Justificar a checagem de manipulação como etapa
metodológica necessária neste trabalho (seção \ref{sec:bfi44}) exige
recorrer a uma linha de pesquisa
distinta, ainda recente, sobre o uso de LLMs como sujeitos simulados em
ciências comportamentais.

\subsection{LLMs como sujeitos simulados}

\citeonline{argyle2023} cunharam o termo \textit{silicon sampling} para
descrever essa abordagem: ao condicionar um LLM (GPT-3) com descrições
sociodemográficas de participantes reais, principalmente do
\textit{American National Election Studies}, os autores obtiveram
distribuições de resposta fortemente correlacionadas com as das amostras
humanas originais. Esse resultado inaugural sustentou a expectativa de
que modelos de linguagem, adequadamente instruídos, poderiam funcionar
como substitutos de baixo custo para respondentes humanos em pesquisas de
opinião e comportamento. \citeonline{ong2024} situa esse uso dentro de
uma tipologia mais ampla, que divide a pesquisa com LLMs em três
paradigmas: \textit{GPT-ology}, LLMs como modelos computacionais e
\textit{silicon sampling} (LLMs simulando humanos). O autor destaca que
um mesmo experimento pode levar a inferências diferentes conforme a
pergunta de pesquisa e que, no caso do \textit{silicon sampling}, a
suposição de que as respostas do modelo reproduzem as respostas humanas
precisa ser testada, e não assumida.

Estudos subsequentes qualificam esse otimismo inicial. \citeonline{park2024}
reproduziram 14 estudos do projeto de replicação \textit{Many Labs 2}
usando GPT-3.5 e, entre os oito estudos que puderam ser analisados,
replicaram apenas 37,5\% dos resultados originais; mais relevante para
este trabalho, os autores documentaram que o modelo respondia a perguntas
sobre orientação política, preferências econômicas, julgamento e
filosofia moral com variação nula ou quase nula entre execuções,
convergindo para uma mesma resposta supostamente ``correta'', padrão que
chamam de efeito da ``resposta correta'' (\textit{correct answer
effect}) e que motiva a preocupação, expressa no título do artigo, com
uma diversidade de pensamento diminuída. \citeonline{petrov2024},
especificamente sobre instrumentos de personalidade, testaram GPT-3.5 e
GPT-4 respondendo a questionários padronizados de traços, incluindo o
Big Five, sob diferentes personas: no GPT-4, mas não no GPT-3.5,
descrições genéricas de pessoa produziram respostas com propriedades
psicométricas promissoras, ainda que imperfeitas, próximas de normas
humanas; já perfis demográficos específicos produziram propriedades
psicométricas ruins nos dois modelos. Os autores concluem que,
atualmente, quando um LLM é instruído a simular personas definidas por
perfis demográficos, suas respostas são sinais pobres dos traços latentes
que supostamente representariam.

Os achados de \citeonline{park2024} e \citeonline{petrov2024}, tomados em
conjunto, sustentam diretamente a necessidade da checagem de manipulação
adotada neste trabalho: se a capacidade de um LLM de incorporar uma
personalidade induzida não pode ser assumida \textit{a priori}, tampouco
pode ser assumida a validade de qualquer relação estatística observada
entre essa personalidade e uma variável dependente, sem antes confirmar
que a indução de fato ocorreu. A checagem de manipulação é, portanto,
uma condição para que os resultados do Capítulo \ref{cap:resultados}
sejam interpretáveis. Os três estudos empíricos citados acima
\cite{argyle2023,park2024,petrov2024} testam modelos da família GPT
(OpenAI), e não os modelos abertos usados neste trabalho (seções
\ref{sec:selecao-modelo} e \ref{sec:metodo-robustez}); os achados servem
como motivação e paralelo conceitual, não como previsão direta do
comportamento dos modelos efetivamente usados aqui. O próprio
\textit{pipeline} deste trabalho também documentou vieses de resposta
dos modelos (viés de centralidade e maior dificuldade com itens de
pontuação reversa em modelos menores), discutidos nas seções
\ref{sec:limitacoes-pipeline} e \ref{sec:vieses}.

\subsection{Personalidade induzida e risco em LLMs}
\label{sec:ref-llm-risco}

Dois estudos recentes aproximam essa linha de pesquisa do tema deste
trabalho, ao investigar especificamente a relação entre personalidade
induzida em LLMs e comportamento de risco financeiro.

\citeonline{borman2024} construíram personas de LLM a partir do modelo
Big Five, atribuindo a cada um dos cinco traços um nível Baixo, Médio ou
Alto e gerando as 243 combinações possíveis desses níveis. As personas
foram testadas com ChatGPT 3.5 e ChatGPT 4.0 (os autores reportam os
resultados do segundo), e a personalidade induzida foi validada por meio
de um teste de personalidade do modelo dos cinco fatores, uma checagem
de manipulação análoga à deste trabalho. As personas responderam a dois
instrumentos: um questionário comportamental curto, de nove perguntas
elaboradas principalmente a partir de estudos anteriores com humanos, e
uma tarefa simulada de investimento, criada pelos autores, em que cada
persona escolhia uma entre cinco empresas com retorno e fator de risco
informados, podendo optar por pesquisar sobre elas antes de decidir. A
relação entre traços e comportamentos foi estimada por regressão OLS nos
dois casos. No apetite ao risco, o questionário reproduziu a direção
esperada pela literatura humana em três traços (Abertura à Experiência,
Extroversão e Neuroticismo), mas não em Consciência e Amabilidade; já a
tarefa simulada reproduziu a direção esperada nos cinco traços. De forma
mais geral, os autores concluem que as personas representaram o
comportamento humano melhor na tarefa simulada do que no questionário.

\citeonline{hartley2025} estudaram a relação entre personalidade e
propensão ao risco em diversos LLMs comerciais (entre eles GPT-4o,
GPT-4o mini, GPT-4 Turbo, Claude 3 Sonnet, Claude 3 Haiku, Gemini 1.5 Pro
e Gemini 1.5 Flash), usando a Teoria da Perspectiva Cumulativa: o risco
é medido pelos parâmetros que melhor explicam o equivalente certo
informado pelo modelo para cada loteria, isto é, o valor garantido que
ele aceitaria no lugar de apostar. Diferentemente de
\citeonline{borman2024} e do Experimento A deste trabalho, a indução é
feita um traço por vez, com os demais mantidos constantes, por meio de
adjetivos bipolares em uma escala de intensidade de 1 a 9. Os autores
relatam três resultados relevantes para este trabalho. Primeiro, no
GPT-4o, Abertura à Experiência foi o único traço significativamente
correlacionado com os parâmetros de sensibilidade ao risco, em linha com
a literatura humana citada por eles; a comparação entre os cinco traços
é relatada pelos autores apenas para o GPT-4o, e nos demais modelos a
análise relatada concentra-se em Abertura à Experiência. Segundo, a
relação entre
personalidade e risco não foi uniforme entre modelos: as variantes
menores (GPT-4o mini e Claude 3 Haiku) não reproduziram o padrão das
variantes maiores, e alguns modelos apresentaram relações na direção
oposta à humana. Terceiro, sem intervenção de personalidade, a maioria
dos modelos se comportou de forma aproximadamente neutra ao risco.

Os dois estudos se diferenciam deste trabalho em pontos centrais. Usam
modelos proprietários comerciais, e não modelos de pesos abertos, como os
usados neste trabalho. Além disso, medem o risco por uma tarefa de escolha
entre empresas, em \citeonline{borman2024}, ou pelos equivalentes certos
de loterias informados pelo modelo, em \citeonline{hartley2025} e, no primeiro caso, também por um
questionário comportamental; nenhum deles usa um questionário de
\textit{suitability} como o IPA. Ainda assim, oferecem dois paralelos
diretos que são retomados no Capítulo \ref{cap:discussao}: o desenho de
indução de \citeonline{borman2024} (os cinco traços simultâneos em três
níveis) é essencialmente o mesmo do Experimento A (Capítulo
\ref{cap:metodologia}), e o papel de Abertura à Experiência como traço
mais associado ao risco aparece também em \citeonline{hartley2025}, no
GPT-4o.
